\chapter{2021年度}
\label{ch:2021}

\section{第1期・専門基礎科目（3題必修）}

\begin{problem}{第1問｜Jordan 標準形}{2021-1-1}
$a\in\mathbb C$ とし
\[
A=\begin{pmatrix}a-1&a&0\\-a&-a-1&0\\-1&1&-1\end{pmatrix}.
\]
固有多項式，固有値と固有空間の次元，Jordan 標準形を求めよ。
\end{problem}
\begin{memo*}
\textbf{解答の見通し。}
行列 $A$ の対角成分には全て $-1$ が含まれているので、$N=A+I_3$ とおいて
$-I_3$ の部分を取り除く。実際に掛け算すると $N^3=O$ になる。
このように何乗かすると零行列になる行列を冪零行列という。
$Nx=\mu x$ なら $N^3x=\mu^3x=0$ なので、$N$ の固有値は0だけである。
従って $A=N-I_3$ の固有値は $-1$ だけになる。

Jordan ブロック $J_k(-1)$ に $I_k$ を足すと、上の対角線に1が並ぶ冪零ブロックになる。
その $k$ 乗で初めて0になるため、$N^2$ と $N^3$ を調べれば最大ブロックの大きさが分かる。
$a=0$ のときだけ $N^2$ が0になるので場合分けする。
\end{memo*}
\begin{proof}[解答]
$N=A+I_3$ とおくと、$A$ の対角成分に $1$ を足して
\[
N=\begin{pmatrix}a&a&0\\-a&-a&0\\-1&1&0\end{pmatrix}
\]
である。この行列の冪を計算する。第1行 $(a,a,0)$ と $N$ の第1列 $(a,-a,-1)^{\mathsf T}$ の積は
$a^2-a^2+0=0$ であり、第2列 $(a,-a,1)^{\mathsf T}$ との積も $a^2-a^2+0=0$、
第3列は零ベクトルなので0である。第2行 $(-a,-a,0)$ は第1行の $-1$ 倍だから、
その行も全て0になる。第3行 $(-1,1,0)$ については
\[
(-1)a+1\cdot(-a)+0=-2a,\qquad
(-1)a+1\cdot(-a)+0=-2a
\]
となる。従って
\[
N^2=\begin{pmatrix}0&0&0\\0&0&0\\-2a&-2a&0\end{pmatrix}.
\]
さらに $N^3=N^2N$ の第3行は $(-2a,-2a,0)$ と $N$ の各列の積であり、
第1列では $-2a\cdot a+(-2a)(-a)=-2a^2+2a^2=0$、第2列でも同様に0、第3列は0である。
他の行はすでに0なので
\[
N^3=O
\]
である。

$Nx=\mu x$ を満たす $x\ne0$ があれば
\[
0=N^3x=\mu^3x
\]
だから $\mu^3=0$、すなわち $\mu=0$ である。従って $N$ の固有値は0だけであり、
3次の特性多項式は根を重複度込みで全て並べたものだから $|\mu I_3-N|=\mu^3$ となる。
$A=N-I_3$ なので $\lambda I_3-A=(\lambda+1)I_3-N$ であり、上の式へ $\mu=\lambda+1$ を代入して
\[
|\lambda I_3-A|=(\lambda+1)^3.
\]
従って固有値は $-1$ のみ（重複度3）である。

$a\ne0$ の場合を見る。このとき $N^2\ne O$ かつ $N^3=O$ である。
大きさ $k$ の冪零 Jordan ブロックは $k$ 乗で初めて0になるので、
最大の Jordan ブロックの大きさは3である。
全体の大きさも3だから、標準形はブロック一つの $J_3(-1)$ に限る。
固有空間の次元も直接分かる。$Nx=0$ は
\[
ax_1+ax_2=0,\qquad -ax_1-ax_2=0,\qquad -x_1+x_2=0
\]
であり、$a\ne0$ より第1式は $x_1+x_2=0$、第3式は $x_1=x_2$ だから $x_1=x_2=0$、
$x_3$ は自由である。従って
\[
\ker(A+I_3)=\ker N=\left\langle(0,0,1)^{\mathsf T}\right\rangle,
\qquad \dim\ker(A+I_3)=1.
\]

$a=0$ なら
\[
N=\begin{pmatrix}0&0&0\\0&0&0\\-1&1&0\end{pmatrix},
\]
従って $N\ne O$, $N^2=O$ である。
$x=(x_1,x_2,x_3)^{\mathsf T}$ とすると
\[
Nx=\begin{pmatrix}0\\0\\-x_1+x_2\end{pmatrix}=0
\quad\Longleftrightarrow\quad x_2=x_1.
\]
従って全ての解は
\[
x=x_1\begin{pmatrix}1\\1\\0\end{pmatrix}
 +x_3\begin{pmatrix}0\\0\\1\end{pmatrix}
\]
と表され、自由に選べる数が $x_1,x_3$ の二つあるので
$\dim\ker N=2$ である。
$N\ne O$ なので大きさ2のブロックが少なくとも一つ必要である。
大きさ $k$ の冪零 Jordan ブロック
\[
J_k(0)=\begin{pmatrix}
0&1&&0\\
&0&\ddots&\\
&&\ddots&1\\
0&&&0
\end{pmatrix}
\]
では $J_k(0)y=0$ の解は $y=(y_1,0,\ldots,0)^{\mathsf T}$ であり、
各ブロックが核の自由な数を一つずつ作る。従って
$\dim\ker N=2$ は Jordan ブロックが2個あることを表す。
全体の大きさは3だから、ブロックの大きさは2と1になり、標準形は
$J_2(-1)\oplus(-1)$ となる。
\end{proof}

\begin{problem}{第2問｜二変数関数}{2021-1-2}
\[
f(x,y)=\begin{cases}\dfrac{\sin(x^2y)}{x^2+y^2},&(x,y)\ne(0,0),\\0,&(x,y)=(0,0).
\end{cases}
\]
\begin{enumerate}
  \item[(1)] 極座標で $|f(x,y)|\le r$ および $|f(x,y)|\le r^{-2}$ を示せ。
  \item[(2)] 原点で連続であることを示せ。
  \item[(3)] $p>1$，$B(R)=\{x^2+y^2\le R^2\}$ とし $I_p(R)=\iint_{B(R)}|f|^p\,\d x\d y$ とおく。$\lim_{R\to\infty}I_p(R)$ の存在を示せ。
  \item[(4)] $\partial f/\partial y(x,0)$ を求め，$f$ が $C^1$ 級か判定せよ。
\end{enumerate}
\end{problem}
\begin{memo*}
\textbf{解答の見通し。}
極座標の半径 $r=\sqrt{x^2+y^2}$ は、点 $(x,y)$ と原点の距離であるため
$|x|\le r$, $|y|\le r$ が成り立つ。
$|\sin u|\le|u|$ を使うと $|f|\le r$ となり、原点へ近づくと0へ近づくことが分かる。
一方、$|\sin u|\le1$ を使うと $|f|\le r^{-2}$ となり、遠方での減衰が分かる。
半径1の内側と外側で二つの評価を使い分け、面積要素
$\d x\d y=r\,\d r\d\theta$ を掛けた一変数積分が収束するか確認する。
最後の偏導関数は公式を先に使わず、定義の差商を直接計算する。
\end{memo*}
\begin{proof}[解答]
\begin{enumerate}
  \item[(1)] $r=\sqrt{x^2+y^2}>0$ とおくと $|x|\le r$, $|y|\le r$ である。
  $|\sin u|\le|u|$ を $u=x^2y$ に使えば
  \[
  |f(x,y)|
  \le\frac{|x|^2|y|}{x^2+y^2}
  \le\frac{r^3}{r^2}=r.
  \]
  また $|\sin u|\le1$ より
  \[
  |f(x,y)|\le\frac1{x^2+y^2}=\frac1{r^2}.
  \]
  \item[(2)] 任意の $\varepsilon>0$ に対して $\delta=\varepsilon$ と取る。
  $0<\sqrt{x^2+y^2}=r<\delta$ なら (1) より
  \[
  |f(x,y)-f(0,0)|=|f(x,y)|\le r<\varepsilon.
  \]
  従って $f$ は原点で連続である。
  \item[(3)] $r\le1$ では $|f|^p\le r^p$、$r\ge1$ では
  $|f|^p\le r^{-2p}$ である。極座標の面積要素は
  $\d x\d y=r\,\d r\d\theta$ だから
  \begin{align*}
  \iint_{r\le1}|f|^p\,\d x\d y
  &\le\int_0^{2\pi}\int_0^1r^p\cdot r\,\d r\d\theta\\
  &=2\pi\int_0^1r^{p+1}\,\d r
  =2\pi\left[\frac{r^{p+2}}{p+2}\right]_0^1
  =\frac{2\pi}{p+2}<\infty,
  \end{align*}
  \begin{align*}
  \iint_{r\ge1}|f|^p\,\d x\d y
  &\le\int_0^{2\pi}\int_1^\infty r^{-2p}\cdot r\,\d r\d\theta\\
  &=2\pi\int_1^\infty r^{1-2p}\,\d r
  =2\pi\left[\frac{r^{2-2p}}{2-2p}\right]_1^\infty
  =\frac{2\pi}{2p-2}<\infty.
  \end{align*}
  最後の広義積分では、上端で $r^{2-2p}\to0$ となるために指数の条件
  $2-2p<0$、すなわち $p>1$ が必要であり、仮定はこれを満たしている。
  従って $|f|^p\in L^1(\R^2)$。また
  $R$ を増やすと円板 $B(R)$ は広がるので、非負関数
  $|f|^p\mathbf1_{B(R)}$ は各点で $|f|^p$ へ単調に増加する。従って単調収束定理より
  \[
  \lim_{R\to\infty}I_p(R)=\iint_{\R^2}|f|^p\,\d x\d y<\infty.
  \]
  \item[(4)] 偏微分を計算すると
  \[
  \frac{\partial f}{\partial y}(x,0)=\begin{cases}1,&x\ne0,\\0,&x=0.
  \end{cases}
  \]
  実際、$x\ne0$ では偏微分の定義から
  \begin{align*}
  \frac{\partial f}{\partial y}(x,0)
  &=\lim_{h\to0}\frac{f(x,h)-f(x,0)}h\\
  &=\lim_{h\to0}\frac{\sin(x^2h)}{h(x^2+h^2)}\\
  &=\lim_{h\to0}
    \frac{\sin(x^2h)}{x^2h}\frac{x^2}{x^2+h^2}=1.
  \end{align*}
  最後の等号では $\sin u/u\to1$（$u=x^2h\to0$）と
  $x^2/(x^2+h^2)\to1$（$x\ne0$ を固定）を使った。
  $x=0$ では、全ての $h$ について $f(0,h)=\sin0/h^2=0$ かつ $f(0,0)=0$ なので、
  同じ差商は0であり $\partial f/\partial y(0,0)=0$ である。

  $C^1$ 級とは、全ての1階偏導関数が存在して連続になることである。
  いま $x\to0$（$x\ne0$）とすると
  \[
  \frac{\partial f}{\partial y}(x,0)=1
  \longrightarrow1
  \ne0=\frac{\partial f}{\partial y}(0,0)
  \]
  だから、$\partial f/\partial y$ は原点で不連続である。従って $f$ は $C^1$ 級ではない。
\end{enumerate}
\end{proof}

\begin{problem}{第3問｜連続像}{2021-1-3}
\begin{enumerate}
  \item[(1)] 集合と写像 $f:X\to Y$, $g:Y\to Z$ で，$f$ は全射でなく，$g$ は単射でなく，$g\circ f$ は全単射となる例を挙げよ。
  \item[(2)] 連続写像による連結集合の像は連結か。
  \item[(3)] 連続写像によるコンパクト集合の像はコンパクトか。
\end{enumerate}
\end{problem}
\begin{memo*}
\textbf{解答の見通し。}
(1) で、全射は「終域の全ての元が像に現れる」、単射は
「異なる二点が同じ値へ写らない」、全単射は両方を満たすことを意味する。
一元集合と二元集合を使えば、各条件を元を列挙して確認できる。
(2) の連結とは、空でない二つの相対開集合へ分割できないことである。
像が分割できると仮定し、その二集合の逆像を取ると元の集合まで分割される。
(3) のコンパクト性は、任意の開被覆から有限部分被覆を選べることとして使う。
像の開被覆を逆像へ戻して定義域を覆い、選んだ有限個を再び像へ戻す。
\end{memo*}
\begin{proof}[解答]
\begin{enumerate}
  \item[(1)] $X=Z=\{0\}$、$Y=\{0,1\}$ とし
  \[
  f(0)=0,\qquad g(0)=g(1)=0
  \]
  と定める。$1\notin f(X)$ なので $f$ は全射でなく、$g(0)=g(1)$ なので $g$ は単射でない。
  一方、$g\circ f:X\to Z$ は一元集合から一元集合への写像なので全単射である。

  \item[(2)] 連結である。連結集合 $C$ と連続写像 $h:C\to Y$ を取り、$h(C)$ が連結でないと仮定する。
  このとき $h(C)$ の相対位相で開な非空集合 $U,V$ が存在して
  \[
  h(C)=U\sqcup V
  \]
  と書ける。$U=h(C)\cap U'$, $V=h(C)\cap V'$ となる $Y$ の開集合 $U',V'$ を取れば
  \[
  C=h^{-1}(U')\sqcup h^{-1}(V')
  \]
  となる。右辺の二集合は連続性により $C$ で開であり、$U,V$ が非空なので、ともに非空である。
  これは $C$ の連結性に矛盾する。従って $h(C)$ は連結である。

  \item[(3)] コンパクトである。コンパクト集合 $K$ と連続写像 $h:K\to Y$ を取り、
  $h(K)$ の開被覆を $\{U_\lambda\}_{\lambda\in\Lambda}$ とする。このとき
  \[
  K\subset\bigcup_{\lambda\in\Lambda}h^{-1}(U_\lambda)
  \]
  であり、各 $h^{-1}(U_\lambda)$ は $K$ で開である。$K$ のコンパクト性より有限個
  $\lambda_1,\ldots,\lambda_m$ を選んで
  \[
  K\subset\bigcup_{j=1}^m h^{-1}(U_{\lambda_j})
  \]
  とできる。両辺を $h$ で写せば
  \[
  h(K)\subset\bigcup_{j=1}^m U_{\lambda_j}
  \]
  となるので、$h(K)$ はコンパクトである。
\end{enumerate}
\end{proof}

\section{第1期・専門科目（4題から2題選択）}

\begin{problem}{第1問｜環とイデアル}{2021-s-1}
$\Phi:\mathbb C[X,Y,Z]\to\mathbb C[t]$ を
$\Phi(f)=f(t^2,t^3,t^4)$ で定め，$I=\ker\Phi$，$R=\mathbb C[X,Y,Z]/I$ とする。
\begin{enumerate}
  \item[(1)] $I$ は素イデアルか。
  \item[(2)] $I$ の生成元を求めよ。
  \item[(3)] $\bar X,\bar Y,\bar Z$ で生成される $R$ のイデアル $\mathfrak m$ が極大であることを示せ。
  \item[(4)] $\mathfrak m$ が2元で生成されることを示せ。
  \item[(5)] $\mathfrak m$ が単項イデアルでないことを示せ。
\end{enumerate}
\end{problem}
\begin{memo*}
\textbf{解答の見通し。}
準同型の像を $\mathbb C[t^2,t^3]$ と同定して商環を具体化する。
核ではまず $Z=X^2$ を消去し、次に $Y^2=X^3$ で割った余りを偶数次数・奇数次数に分ける。
極大性は原点での値写像、生成元の最小個数は
$\mathfrak m/\mathfrak m^2$ のベクトル空間次元で判定する。
\end{memo*}
\begin{proof}[解答]
\begin{enumerate}
  \item[(1)] $t^4=(t^2)^2$ だから
  \[
  \operatorname{Im}\Phi=\mathbb C[t^2,t^3,t^4]
  =\mathbb C[t^2,t^3].
  \]
  第1同型定理より
  \[
  R=\mathbb C[X,Y,Z]/\ker\Phi
  \cong\mathbb C[t^2,t^3]\subset\mathbb C[t].
  \]
  $\mathbb C[t]$ は整域なので、その部分環 $\mathbb C[t^2,t^3]$ も整域である。
  従って $I=\ker\Phi$ は素イデアルである。
  \item[(2)] 生成元は
  \[
  I=(Z-X^2,\ Y^2-X^3).
  \]
  実際、$\Phi(Z-X^2)=t^4-(t^2)^2=0$、
  $\Phi(Y^2-X^3)=t^6-(t^2)^3=0$ だから、右辺を $J$ とおけば $J\subset I$。

  逆に $f\in I$ とする。$Z-X^2$ を用いて $Z$ を消去すれば、ある
  $g(X,Y)\in\mathbb C[X,Y]$ により
  \[
  f\equiv g(X,Y)\pmod{(Z-X^2)}
  \]
  と書ける。さらに $g$ を $Y$ の多項式と見て、モニック多項式
  $Y^2-X^3$ で割ると
  \[
  g(X,Y)=q(X,Y)(Y^2-X^3)+a(X)Y+b(X)
  \]
  となる。$f\in I$ より
  \[
  0=\Phi(f)=a(t^2)t^3+b(t^2).
  \]
  $a(t^2)t^3$ は奇数次数の項だけ、$b(t^2)$ は偶数次数の項だけからなるので、
  この多項式恒等式から $a=b=0$。従って $f\in J$ であり、$I=J$ を得る。

  \item[(3)] $R\cong\mathbb C[t^2,t^3]$ のもとで
  $\mathfrak m=(t^2,t^3)$ に対応する。定数項を取る準同型
  \[
  \mathbb C[t^2,t^3]\longrightarrow\mathbb C,\qquad h(t)\longmapsto h(0)
  \]
  は全射で、核は $(t^2,t^3)$ である。従って
  $R/\mathfrak m\cong\mathbb C$ は体だから $\mathfrak m$ は極大である。

  \item[(4)] $\bar Z=\bar X^2$ だから
  \[
  \mathfrak m=(\bar X,\bar Y,\bar Z)
  =(\bar X,\bar Y,\bar X^2)=(\bar X,\bar Y).
  \]

  \item[(5)] $R\cong\mathbb C[t^2,t^3]$ で考えると
  \[
  \mathfrak m=(t^2,t^3),\qquad
  \mathfrak m^2=(t^4,t^5,t^6).
  \]
  従って $\mathfrak m/\mathfrak m^2$ では $t^2,t^3$ の剰余類が一次独立であり、
  これらが全体を生成するから
  \[
  \dim_{R/\mathfrak m}(\mathfrak m/\mathfrak m^2)
  =\dim_{\mathbb C}(\mathfrak m/\mathfrak m^2)=2.
  \]
  もし $\mathfrak m=(h)$ が単項なら、$\mathfrak m/\mathfrak m^2$ は
  $h+\mathfrak m^2$ 一つで生成される $R/\mathfrak m$ ベクトル空間なので次元は高々1となり、矛盾する。
\end{enumerate}
\end{proof}

\begin{problem}{第2問｜曲率・捩率}{2021-s-2}
$p(t)=(\cos t,\sin t,\cosh t)$ と円柱面 $S=\{x^2+y^2=1\}$ を考える。
\begin{enumerate}
  \item[(1)] $p$ の曲率と捩率を求めよ。
  \item[(2)] $S$ の主方向と $p'(t)$ の角度が $\pi/4$ となる $t$ を全て求めよ。
  \item[(3)] $p$ の $S$ における測地的曲率ベクトルを求めよ。
\end{enumerate}
\end{problem}
\begin{memo*}
\textbf{解答の見通し。}
(1) は $p',p'',p'''$ を曲率・捩率の公式へ代入する。
(2) では円柱の二つの主方向へ $p'$ を分解し、内積から角度条件を作る。
(3) は単位接ベクトルを弧長で微分して曲率ベクトルを求め、
そこから円柱の法線方向成分を引けば測地的曲率ベクトルになる。
\end{memo*}
\begin{proof}[解答]
まず
\[
p'=(-\sin t,\cos t,\sinh t),\quad
p''=(-\cos t,-\sin t,\cosh t),
\]
\[
p'''=(\sin t,-\cos t,\sinh t)
\]
であり
\[
p'\times p''=(\cos t\cosh t+\sin t\sinh t,
\sin t\cosh t-\cos t\sinh t,1).
\]
従って
\[
\norm{p'}=\cosh t,\qquad
\norm{p'\times p''}=\sqrt2\cosh t,\qquad
(p'\times p'')\cdot p'''=2\sinh t.
\]
曲率・捩率の公式に代入すると
\[
\kappa(t)=\frac{\sqrt2}{\cosh^2t},\qquad
\tau(t)=\frac{\sinh t}{\cosh^2t}.
\]
円柱の単位主方向を
\[
e_1=(-\sin t,\cos t,0),\qquad e_2=(0,0,1)
\]
とする。$p'=e_1+\sinh t\,e_2$ かつ $\norm{p'}=\cosh t$ なので、
$p'$ と $e_1$ のなす角が $\pi/4$ である条件は
\[
\frac{|p'\cdot e_1|}{\norm{p'}}=\frac1{\cosh t}=\frac1{\sqrt2}.
\]
$e_2$ との角についても $|\sinh t|/\cosh t=1/\sqrt2$ となり、同じ条件
$\cosh t=\sqrt2$ を与える。従って
$t=\pm\operatorname{arcosh}\sqrt2=\pm\log(1+\sqrt2)$。

単位接ベクトルは
\[
T=\frac{p'}{\norm{p'}}
=(-\operatorname{sech}t\sin t,
  \operatorname{sech}t\cos t,\tanh t).
\]
$\d s/\d t=\norm{p'}=\cosh t$ より $\d/\d s=\operatorname{sech}t\,\d/\d t$。
従って成分ごとに微分すると
\[
\frac{\d T}{\d s}
=\operatorname{sech}t\frac{\d T}{\d t}
\]
\[
=\bigl(
\operatorname{sech}^2t\tanh t\sin t-\operatorname{sech}^2t\cos t,
-\operatorname{sech}^2t\tanh t\cos t-\operatorname{sech}^2t\sin t,
\operatorname{sech}^3t
\bigr).
\]
円柱の単位法線は $n=(\cos t,\sin t,0)$ であり
\begin{align*}
\frac{\d T}{\d s}\cdot n
&=\operatorname{sech}^2t\tanh t\sin t\cos t
-\operatorname{sech}^2t\cos^2t\\
&\quad-\operatorname{sech}^2t\tanh t\cos t\sin t
-\operatorname{sech}^2t\sin^2t\\
&=-\operatorname{sech}^2t.
\end{align*}
従って法線成分を引けば
\[
\boldsymbol{k}_g
=\frac{\d T}{\d s}
-\left(\frac{\d T}{\d s}\cdot n\right)n
=\operatorname{sech}^2t\,
(\tanh t\sin t,-\tanh t\cos t,\operatorname{sech}t).
\]
\end{proof}

\begin{problem}{第3問｜収束べき級数}{2021-s-3}
$\varphi(z)=\sum_{n=0}^\infty a_nz^n$ は収束べき級数で $a_0\ne0$ とする。
\begin{enumerate}
  \item[(1)] $1/\varphi(z)$ が $z=0$ で収束べき級数に展開できることを示せ。
  \item[(2)] その収束半径が $\varphi$ の収束半径より真に小さくなることはあるか。
  \item[(3)] $1/\varphi(z)=\sum b_nz^n$ とするとき $b_0,b_1,b_2$ を求めよ。
  \item[(4)] $n=0,1,2,3$ に対し $\displaystyle\int_{|z|=1}\frac{\d z}{z^n\sin^2z}$ を求めよ。
\end{enumerate}
\end{problem}
\begin{memo*}
\textbf{解答の見通し。}
$\varphi(0)\ne0$ なら逆数は原点近くで正則だが、その収束半径は
$\varphi$ 自身の収束半径ではなく、逆数が最初に出会う特異点で決まる。
係数は積が1になる条件から再帰的に求める。
積分では $1/\sin^2z$ を必要な次数まで Laurent 展開し、$z^{-1}$ の係数だけを読む。
\end{memo*}
\begin{proof}[解答]
\begin{enumerate}
\item[(1)] $\varphi(0)\ne0$ なので、連続性により0の十分小さい近傍で $\varphi(z)\ne0$。
従って $1/\varphi$ はその近傍で正則であり、Taylor 展開できる。
\item[(2)] ある。例えば $\varphi(z)=1-2z$ は整関数だが
$1/\varphi(z)=\sum_{n=0}^\infty(2z)^n$ の収束半径は $1/2$。
\item[(3)] 恒等式
\[
(a_0+a_1z+a_2z^2+\cdots)(b_0+b_1z+b_2z^2+\cdots)=1
\]
で $z^0,z^1,z^2$ の係数を比較すると
\[
a_0b_0=1,\quad a_0b_1+a_1b_0=0,\quad
a_0b_2+a_1b_1+a_2b_0=0.
\]
これを順に解いて
\[
b_0=\frac1{a_0},\quad b_1=-\frac{a_1}{a_0^2},\quad
b_2=\frac{a_1^2-a_0a_2}{a_0^3}.
\]
\item[(4)]
\[
\sin z=z-\frac{z^3}{6}+\frac{z^5}{120}+O(z^7)
\]
を二乗すると
\begin{align*}
\sin^2z
&=z^2-\frac13z^4
+\left(\frac1{36}+\frac1{60}\right)z^6+O(z^8)\\
&=z^2\left(1-\frac13z^2+\frac{2}{45}z^4+O(z^6)\right).
\end{align*}
$u=-z^2/3+2z^4/45+O(z^6)$ とおき
$(1+u)^{-1}=1-u+u^2+O(u^3)$ を使うと
\begin{align*}
\left(1-\frac13z^2+\frac{2}{45}z^4+O(z^6)\right)^{-1}
&=1+\frac13z^2
+\left(\frac19-\frac{2}{45}\right)z^4+O(z^6)\\
&=1+\frac13z^2+\frac1{15}z^4+O(z^6).
\end{align*}
従って
\[
\frac1{\sin^2z}=\frac1{z^2}+\frac13+\frac{z^2}{15}+O(z^4).
\]
従って
\[
\frac1{z^n\sin^2z}
=z^{-n-2}+\frac13z^{-n}+\frac1{15}z^{2-n}+O(z^{4-n}).
\]
$z^{-1}$ の係数、すなわち留数を $n=0,1,2,3$ の順に読むと
$0,1/3,0,1/15$。留数定理から積分値は順に
\[
0,\qquad \frac{2\pi i}{3},\qquad0,\qquad\frac{2\pi i}{15}.
\]
\end{enumerate}
\end{proof}

\begin{problem}{第4問｜測度収束}{2021-s-4}
$E\subset\R^d$ は可測で $|E|<\infty$，$f_n$ は $E$ 上可測とする。$A_{n,\varepsilon}=\{x\in E\mid|f_n(x)|\ge\varepsilon\}$ とおく。
\begin{enumerate}
  \item[(1)] $g\in L^1(E)$，$F=\{|g|\ge1\}$ なら $|F|\le\int_E|g|$ を示せ。
  \item[(2)] $\int_E|f_n|\to0$ なら全ての $\varepsilon>0$ で $|A_{n,\varepsilon}|\to0$ を示せ。
  \item[(3)] $f_n\to0$ がほとんど至る所なら $\displaystyle\int_E\frac{|f_n|}{1+|f_n|}\to0$ を示せ。
  \item[(4)] 全ての $\varepsilon>0$ で $|A_{n,\varepsilon}|\to0$ なら，ほとんど至る所で0に収束する部分列があることを示せ。
\end{enumerate}
\end{problem}
\begin{memo*}
\textbf{解答の見通し。}
(1), (2) は集合の定義関数を積分して Chebyshev 型の評価を得る。
(3) は被積分関数が1以下なので、有限測度という仮定が優関数を与える。
(4) は閾値を $1/k$、例外集合の測度を $2^{-k}$ とする部分列を選び、
例外集合へ無限回入る点が零集合であることを使う。
\end{memo*}
\begin{proof}[解答]
\begin{enumerate}
\item[(1)] $F=\{|g|\ge1\}$ 上では $1\le|g|$ だから、$E$ 上で
$\mathbf1_F\le|g|$ である。両辺を積分して
\[
|F|=\int_E\mathbf1_F\le\int_E|g|.
\]

\item[(2)] $A_{n,\varepsilon}$ 上では $|f_n|/\varepsilon\ge1$ なので
\[
\mathbf1_{A_{n,\varepsilon}}
\le\frac{|f_n|}{\varepsilon}.
\]
従って
\[
|A_{n,\varepsilon}|
\le\frac1\varepsilon\int_E|f_n|\longrightarrow0.
\]

\item[(3)]
\[
0\le\frac{|f_n(x)|}{1+|f_n(x)|}\le1
\]
であり、$f_n(x)\to0$ となる点では左辺も0へ収束する。
$|E|<\infty$ なので優関数 $\mathbf1_E$ は可積分であり、優収束定理から
\[
\int_E\frac{|f_n|}{1+|f_n|}\longrightarrow0.
\]

\item[(4)] 各 $k$ について $|A_{n,1/k}|\to0$ だから、$n_1<n_2<\cdots$ を
\[
|A_{n_k,1/k}|<2^{-k}
\]
となるよう帰納的に選べる。すると
\[
\sum_{k=1}^\infty|A_{n_k,1/k}|
\le\sum_{k=1}^\infty2^{-k}<\infty.
\]
Borel--Cantelli の補題より、ほとんど全ての $x$ は
$A_{n_k,1/k}$ に有限回しか属さない。従ってそのような $x$ では、十分大きい $k$ に対し
\[
|f_{n_k}(x)|<\frac1k,
\]
ゆえに $f_{n_k}(x)\to0$ である。
\end{enumerate}
\end{proof}

\section{第2期・専門基礎科目（3題必修）}

\begin{problem}{第1問｜等角ベクトルと対称行列}{2021-2-1}
$x_1,\ldots,x_n\in\R^3$ が $|(x_i,x_j)|=r$ $(i\ne j)$，$|(x_i,x_i)|=1$ を満たす。ただし $r\ne1$。
\begin{enumerate}
  \item[(1)] 3次実対称行列全体の次元を求めよ。
  \item[(2)] $\sum_i a_i(x_i,x_j)^2=0$ が全ての $j$ で成り立つなら，全ての $a_i$ が0であることを示せ。
  \item[(3)] 行列 $x_ix_i^{\mathsf T}$ が一次独立であることを示せ。
  \item[(4)] $n\le6$ を示せ。
\end{enumerate}
\end{problem}
\begin{memo*}
\textbf{解答の見通し。}
内積 $(x_i,x_j)=x_i^{\mathsf T}x_j$ は二つのベクトルの角度を測る数である。
目標の $n\le6$ を得るには、$n$ 個の対象を6次元の空間へ入れて一次独立だと示せばよい。
そこで各ベクトルから $3\times3$ 対称行列 $X_i=x_ix_i^{\mathsf T}$ を作る。
行列の一次関係 $\sum_i a_iX_i=O$ に $X_j$ との Frobenius 内積を取ると、
\[
\operatorname{tr}(X_iX_j)=(x_i^{\mathsf T}x_j)^2
\]
となり、問題文 (2) のスカラー方程式へ変わる。
従って (2) から $X_i$ が一次独立だと分かり、その個数は対称行列空間の次元6以下になる。
\end{memo*}
\begin{proof}[解答]
(1) 3次実対称行列は $S^{\mathsf T}=S$、すなわち $s_{ji}=s_{ij}$ を満たす行列である。
従って対角成分3個と、対角より上の成分3個を決めれば残りも決まり、
\[
S=\begin{pmatrix}a&d&e\\d&b&f\\e&f&c\end{pmatrix}
\]
と一意に書ける。この表示は
\[
S=aE_{11}+bE_{22}+cE_{33}
 +d(E_{12}+E_{21})+e(E_{13}+E_{31})+f(E_{23}+E_{32})
\]
と同じであり、右辺の6個の行列は対称行列全体を生成し、
係数が成分としてそのまま現れるので一次独立でもある。従って次元は6である。

(2) 仮定より $(x_j,x_j)^2=|(x_j,x_j)|^2=1$、$i\ne j$ では $(x_i,x_j)^2=|(x_i,x_j)|^2=r^2$
である。従って $j$ を固定した式の $i=j$ の項と $i\ne j$ の項を分けると
\begin{align*}
0
&=\sum_i a_i(x_i,x_j)^2\\
&=a_j+r^2\sum_{i\ne j}a_i\\
&=(1-r^2)a_j+r^2\sum_i a_i. \tag{$*$}
\end{align*}
最後の等号では $\sum_{i\ne j}a_i=\sum_i a_i-a_j$ を使った。
$r=|(x_i,x_j)|\ge0$ と $r\ne1$ より $r^2\ne1$、すなわち $1-r^2\ne0$ である。
式 ($*$) の右辺第2項は $j$ によらないので、$j,k$ に対する二式の差を取れば
\[
(1-r^2)(a_j-a_k)=0,
\]
従って全ての $a_j$ は同じ値 $a$ である。これを ($*$) の前の行へ代入すると
\[
0=a+(n-1)r^2a=\{1+(n-1)r^2\}a.
\]
$n\ge1$, $r\ge0$ なので係数 $1+(n-1)r^2$ は正であり、$a=0$、
従って全ての $a_i$ は0である。

(3) $\sum_i a_ix_ix_i^{\mathsf T}=O$ と仮定する。行列の Frobenius 内積
$\langle C,D\rangle_F=\operatorname{tr}(C^{\mathsf T}D)$ を取る。
$x_ix_i^{\mathsf T}$ は対称なので $C^{\mathsf T}=C$ として扱えて、各 $j$ について
\begin{align*}
0
&=\left\langle\sum_i a_ix_ix_i^{\mathsf T},x_jx_j^{\mathsf T}\right\rangle_F\\
&=\sum_i a_i\operatorname{tr}(x_ix_i^{\mathsf T}x_jx_j^{\mathsf T})\\
&=\sum_i a_i(x_i^{\mathsf T}x_j)^2.
\end{align*}
最後の等式は
\[
x_ix_i^{\mathsf T}x_jx_j^{\mathsf T}
=x_i(x_i^{\mathsf T}x_j)x_j^{\mathsf T},\qquad
\operatorname{tr}(x_ix_j^{\mathsf T})=x_j^{\mathsf T}x_i
\]
を順に用いたものである。
これは (2) の式なので全ての $a_i$ は0、従って $x_ix_i^{\mathsf T}$ は一次独立である。

(4) これらは6次元の実対称行列空間に属する $n$ 個の一次独立な元なので $n\le6$。
\end{proof}

\begin{problem}{第2問｜Gamma 関数と Beta 関数}{2021-2-2}
$p,q>0$ とし
\[
\Gamma(p)=\int_0^\infty x^{p-1}e^{-x}\,\d x,
\quad B(p,q)=\int_0^1x^{p-1}(1-x)^{q-1}\,\d x.
\]
$\Gamma(p)\Gamma(q)=B(p,q)\Gamma(p+q)$ を，二重積分に (1) $x=st,y=s(1-t)$，(2) $x=(r\cos\theta)^2,y=(r\sin\theta)^2$ を用いる二通りで示せ。
\end{problem}
\begin{memo*}
\textbf{解答の見通し。}
まず二つの一変数積分の積を、第1象限 $x>0,y>0$ 上の二重積分として書く。
被積分関数が非負なので、Tonelli の定理により積分順序を変えても値は変わらない。
変数変換では、面積が何倍になるかを表す Jacobian の絶対値も必ず掛ける。

一つ目では $s=x+y$ が二点の合計、$t=x/(x+y)$ が合計に占める $x$ の割合であり、
$s$ だけの積分と $t$ だけの積分へ分離できる。
二つ目では $x,y$ を平方付き極座標で表し、動径 $r$ と角度 $\theta$ の積分へ分ける。
最後に $u=r^2$ と $t=\cos^2\theta$ を代入して、定義された Gamma・Beta 積分へ戻す。
\end{memo*}
\begin{proof}[解答]
積の中の積分変数を区別して $x,y$ と書く。被積分関数は非負なので、Tonelli の定理から
\[
\Gamma(p)\Gamma(q)
=\int_0^\infty\int_0^\infty x^{p-1}y^{q-1}e^{-(x+y)}\,\d x\d y.
\]

(1) $x=st$, $y=s(1-t)$ とおくと、$s=x+y>0$, $t=x/(x+y)\in(0,1)$ であり、
\[
\frac{\partial(x,y)}{\partial(s,t)}
=\begin{vmatrix}t&s\\1-t&-s\end{vmatrix}=-s,
\qquad
\left|\frac{\partial(x,y)}{\partial(s,t)}\right|=s.
\]
従って
\begin{align*}
\Gamma(p)\Gamma(q)
&=\int_0^\infty\int_0^1
(st)^{p-1}\{s(1-t)\}^{q-1}e^{-s}s\,\d t\d s\\
&=\left(\int_0^\infty s^{p+q-1}e^{-s}\,\d s\right)
\left(\int_0^1t^{p-1}(1-t)^{q-1}\,\d t\right)\\
&=\Gamma(p+q)B(p,q).
\end{align*}

(2) $x=(r\cos\theta)^2$, $y=(r\sin\theta)^2$ とおく。
$x,y>0$ から $r>0$ かつ $0<\theta<\pi/2$ であり、逆に
\[
r=\sqrt{x+y},\qquad \theta=\arctan\sqrt{\frac yx}
\]
と一意に定まるので、この対応は $(0,\infty)\times(0,\pi/2)$ と第1象限の間の全単射である。
また、後で使うため
\[
x+y=r^2(\cos^2\theta+\sin^2\theta)=r^2
\]
に注意しておく。各偏微分は
\[
x_r=2r\cos^2\theta,\quad x_\theta=-2r^2\sin\theta\cos\theta,
\]
\[
y_r=2r\sin^2\theta,\quad y_\theta=2r^2\sin\theta\cos\theta
\]
だから、Jacobian は
\begin{align*}
\left|\frac{\partial(x,y)}{\partial(r,\theta)}\right|
&=\left|x_ry_\theta-x_\theta y_r\right|\\
&=\left|4r^3\sin\theta\cos^3\theta
 +4r^3\sin^3\theta\cos\theta\right|\\
&=4r^3\sin\theta\cos\theta
(\cos^2\theta+\sin^2\theta)\\
&=4r^3\sin\theta\cos\theta.
\end{align*}
最後で絶対値が外れるのは $0<\theta<\pi/2$ で $\sin\theta,\cos\theta>0$ だからである。
被積分関数を新しい変数で書き直すと
\begin{align*}
x^{p-1}y^{q-1}e^{-(x+y)}
&=(r^2\cos^2\theta)^{p-1}(r^2\sin^2\theta)^{q-1}e^{-r^2}\\
&=r^{2p-2}\cos^{2p-2}\theta\cdot
  r^{2q-2}\sin^{2q-2}\theta\cdot e^{-r^2}
\end{align*}
なので、Jacobian を掛けると
\begin{align*}
&x^{p-1}y^{q-1}e^{-(x+y)}
 \left|\frac{\partial(x,y)}{\partial(r,\theta)}\right|\\
&\qquad=4r^{(2p-2)+(2q-2)+3}
 \cos^{(2p-2)+1}\theta\,\sin^{(2q-2)+1}\theta\,e^{-r^2}\\
&\qquad=4r^{2p+2q-1}
 \cos^{2p-1}\theta\,\sin^{2q-1}\theta\,e^{-r^2}
\end{align*}
となる。この式は $r$ だけの因子と $\theta$ だけの因子の積なので、積分は分離して
\begin{align*}
\Gamma(p)\Gamma(q)
&=4\int_0^\infty r^{2p+2q-1}e^{-r^2}\,\d r
\int_0^{\pi/2}\cos^{2p-1}\theta\sin^{2q-1}\theta\,\d\theta.
\end{align*}
第1積分では $u=r^2$, $\d u=2r\,\d r$ とおく。
$r^{2p+2q-1}\,\d r=r^{2p+2q-2}\cdot r\,\d r=u^{p+q-1}\cdot\frac{\d u}2$ だから
\[
\int_0^\infty r^{2p+2q-1}e^{-r^2}\,\d r
=\frac12\int_0^\infty u^{p+q-1}e^{-u}\,\d u
=\frac12\Gamma(p+q).
\]
第2積分では $t=\cos^2\theta$, $\d t=-2\sin\theta\cos\theta\,\d\theta$ とおく。
このとき $\sin^2\theta=1-t$ であり、
$\theta=0,\pi/2$ がそれぞれ $t=1,0$ に対応するので
\begin{align*}
\int_0^{\pi/2}\cos^{2p-1}\theta\sin^{2q-1}\theta\,\d\theta
&=\int_0^{\pi/2}\cos^{2p-2}\theta\sin^{2q-2}\theta
  \cdot\sin\theta\cos\theta\,\d\theta\\
&=\int_1^0t^{p-1}(1-t)^{q-1}\left(-\frac{\d t}2\right)\\
&=\frac12\int_0^1t^{p-1}(1-t)^{q-1}\,\d t
=\frac12B(p,q).
\end{align*}
係数も含めて $4\cdot(1/2)\cdot(1/2)=1$ なので、再び
$\Gamma(p)\Gamma(q)=\Gamma(p+q)B(p,q)$ を得る。
\end{proof}

\begin{problem}{第3問｜距離空間の連続写像}{2021-2-3}
距離空間 $(X,d_X)$ から $(Y,d_Y)$ への写像 $f$ について答えよ。
\begin{enumerate}
  \item[(1)] $f$ が点 $x\in X$ で連続であることを開球を用いて定義せよ。
  \item[(2)] $f$ が連続なら，任意の閉集合 $F\subset Y$ に対し $f^{-1}(F)$ が閉であることを示せ。
  \item[(3)] 任意の $A\subset X$ に対し $f(\overline A)\subset\overline{f(A)}$ を示せ。
\end{enumerate}
\end{problem}
\begin{memo*}
\textbf{解答の見通し。}
(1) の開球 $B_X(x,\delta)$ は、$x$ からの距離が $\delta$ 未満の点全体である。
連続とは、出力側で許す誤差 $\varepsilon$ を先に指定したとき、
入力側の十分小さい誤差 $\delta$ を選べることをいう。量化記号はこの順序で書く。
(2) は $x\notin f^{-1}(F)$ を一つ取り、その周りに逆像の補集合へ含まれる開球を作る。
(3) の $x\in\overline A$ は、$x$ のどんな開球にも $A$ の点があるという意味である。
連続性が与える $\delta$-球からその点を取り、像が任意の $\varepsilon$-球へ入ることを示す。
\end{memo*}
\begin{proof}[解答]
\begin{enumerate}
\item[(1)] $B_X(x,\delta)=\{y\in X\mid d_X(x,y)<\delta\}$ とする。
$f$ が $x$ で連続であるとは
\[
\forall\varepsilon>0\ \exists\delta>0\quad
f(B_X(x,\delta))\subset B_Y(f(x),\varepsilon)
\]
が成り立つこと、すなわち
\[
d_X(x,y)<\delta\quad\Longrightarrow\quad
d_Y(f(x),f(y))<\varepsilon
\]
が成り立つことである。

\item[(2)] $F\subset Y$ を閉集合とすると $Y\setminus F$ は開集合である。
$x\in f^{-1}(Y\setminus F)$ を取ると $f(x)\in Y\setminus F$ である。
補集合は開なので、ある $\varepsilon>0$ に対して
\[
B_Y(f(x),\varepsilon)\subset Y\setminus F
\]
となる。$x$ における連続性から、ある $\delta>0$ が存在して
\[
f(B_X(x,\delta))\subset B_Y(f(x),\varepsilon)
\]
となる。従って $B_X(x,\delta)\subset f^{-1}(Y\setminus F)$ である。
各点 $x$ の周りにこのような開球が取れるので $f^{-1}(Y\setminus F)$ は開であり、
\[
f^{-1}(Y\setminus F)=X\setminus f^{-1}(F)
\]
である。従って $f^{-1}(F)$ の補集合は開であり、$f^{-1}(F)$ は閉集合である。

\item[(3)] $x\in\overline A$ を取り、$f(x)\in\overline{f(A)}$ を示す。
任意の $\varepsilon>0$ に対して、$x$ における連続性より、ある $\delta>0$ が存在して
\[
f(B_X(x,\delta))\subset B_Y(f(x),\varepsilon)
\]
となる。$x\in\overline A$ だから $B_X(x,\delta)\cap A\ne\varnothing$ である。
この共通部分から $a$ を一つ取ると
\[
f(a)\in B_Y(f(x),\varepsilon)\cap f(A)
\]
となる。従って $f(x)$ の任意の開球は $f(A)$ と交わるので
$f(x)\in\overline{f(A)}$ である。以上より
\[
f(\overline A)\subset\overline{f(A)}
\]
を得る。
\end{enumerate}
\end{proof}
